
%(BEGIN_QUESTION)
% Copyright 2011, Tony R. Kuphaldt, released under the Creative Commons Attribution License (v 1.0)
% This means you may do almost anything with this work of mine, so long as you give me proper credit

Les deler av databladet fra Burr-Brown (Texas Instruments) for INA111 høyhastighets instrumentation amplifier, og svar på følgende spørsmål:

\vskip 10pt

Beregn nødvendig gain-motstand ($R_G$) for forsterkning 30. Oppgi også denne gain-verdien i decibel.

\vskip 10pt

Manipuler den oppgitte gain-formelen for å løse for $R_G$ ved vilkårlig ønsket gain.

\vskip 10pt

Oppgi typisk inngangsresistans for forsterkeren, samt typisk biasstrøm. Beskriv hva som skjer dersom det ikke finnes ``return paths'' for biasstrømmene.

\vskip 10pt

Databladet anbefaler {\it clamping diodes} for å beskytte inngangene mot store overspenninger. Studer figur 5 og forklar hvordan disse diodene beskytter forsterkeren. Merk at det faktisk finnes en feil i dette skjemaet -- finner du den?

\vskip 20pt \vbox{\hrule \hbox{\strut \vrule{} {\bf Forslag til sokratisk diskusjon} \vrule} \hrule}

\begin{itemize}
\item{} Identifiser hvilke grunnleggende kretsprinsipper som brukes i hvert analysetrinn. Forklar altså ``hvorfor'' for hvert trinn, ikke bare ``hva''.
\item{} Databladet sier at én bias-returmotstand kan være nok dersom differensiell kildeimpedans er lav. Forklar hvorfor kildemotstanden betyr noe for biasstrømmene.
\item{} Med internskjemaet på første side: beregn $V_O$ gitt $V_{in(+)}$ = +3.5 V, $V_{in(-)}$ = +1.5 V og $R_G$ = 50 k$\Omega$, ved bruk av Ohms lov og Kirchhoffs lover på tre-opamp-kretsen.
\end{itemize}

\underbar{file i00887}
%(END_QUESTION)





%(BEGIN_ANSWER)

$R_G$ = 1.724 k$\Omega$ for gain 30 = 29.54 dB

\vskip 10pt

$$R_G = {50000 \over {A_v - 1}}$$

%(END_ANSWER)





%(BEGIN_NOTES)

(Gain-formel vist i skjemaet på side 1)

$$A_V = 1 + {50000 \over R_G}$$

$$A_V - 1 = {50000 \over R_G}$$

$${1 \over A_V - 1} = {R_G \over 50000}$$

$$R_G = {50000 \over A_V - 1}$$

$$R_G = {50000 \over 30 - 1} = 1.7241 \hbox{ k}\Omega$$

\vskip 10pt

$$\hbox{dB} = 20 \log A_V$$

$$\hbox{dB} = 20 \log 30 = 29.54 \hbox{ dB}$$

Ifølge side 2 er typisk inngangsresistans for INA111 rundt $10^{12}$ ohm, med biasstrømmer rundt 2 pA (pikoampere). Side 8 beskriver behovet for returvei for biasstrømmer til negativ forsyningsskinne. Uten returveier vil inngangene flyte til potensialer utenfor forsterkerens common-mode-område, og utgangen mettes.

\vskip 10pt

Clamping-diodene er orientert slik at inngangssignaler som overskrider forsyningsskinnene ikke slipper inn i full størrelse til forsterkerinngangene. Overspenninger ``klemmes'' dermed til litt utenfor skinnene. Feilen i skjemaet er at nedre forsyningsskinne skal være merket ``V$-$'', men er feilaktig merket ``V+'' som øvre skinne.

%INDEX% Reading assignment: Burr-Brown INA111 instrumentation amplifier datasheet

%(END_NOTES)
